\documentclass[11pt]{article}
\usepackage[margin=1.05in]{geometry}
\usepackage{amsmath,amssymb,amsthm,mathtools}
\usepackage{xcolor}
\usepackage[colorlinks,linkcolor=blue!55!black,citecolor=blue!55!black]{hyperref}

\theoremstyle{plain}
\newtheorem{thm}{Theorem}[section]
\newtheorem{lem}[thm]{Lemma}
\newtheorem{prop}[thm]{Proposition}
\newtheorem{cor}[thm]{Corollary}
\theoremstyle{definition}
\newtheorem{rem}[thm]{Remark}
\newtheorem{defn}[thm]{Definition}

\DeclareMathOperator{\supp}{supp}
\DeclareMathOperator{\Newton}{Newton}
\DeclareMathOperator{\conv}{conv}
\DeclareMathOperator{\srt}{sort}
\newcommand{\N}{\mathbb{N}}
\newcommand{\Z}{\mathbb{Z}}
\newcommand{\R}{\mathbb{R}}
\newcommand{\cyl}{\mathrm{cyl}}
\newcommand{\lhat}{\hat\lambda}
\newcommand{\Perm}{\mathcal{P}}

\title{Saturated Newton polytope and the Newton polytope identification\\
for cylindric skew Schur polynomials\\
\large (Gap 2 of \texttt{2026-09-20-c1-cylindric-M-convexity}, closed)}
\author{Clio}
\date{25 September 2026 (cycle 2)}

\begin{document}
\maketitle

\begin{abstract}
Let $\lambda/\mu\in\cyl^{n,m}$ be a cylindric skew shape, $\ell\ge 1$ with
$\mathcal{T}_\ell(\lambda/\mu)\ne\emptyset$, and let
$W_\ell=\supp\bigl(s^{c}_{\lambda/\mu}(x_1,\dots,x_\ell)\bigr)\subseteq\N^{\ell}$.
Writing $\lhat$ for the sorted weight of the greedy chain and
$\Perm_{\lhat}=\conv(S_\ell\cdot\lhat)$ for the permutahedron, we prove
\[
  \conv(W_\ell)\cap\Z^{\ell}=W_\ell
  \qquad\text{and}\qquad
  \conv(W_\ell)=\Perm_{\lhat},
\]
the two assertions of Theorem 7.1 of~\cite{note1} that were stated there without
proof. Both arguments are written out in full and neither cites Rado's theorem.
The first needs only the equality of the sorted and subset descriptions of the
weight set, which we isolate as Lemma~\ref{lem:JJ} because it is load-bearing in
two independent places. The second is the classical Hardy--Littlewood--P\'olya
argument, whose engine --- the Robin Hood step --- is Lemma 4.1 of~\cite{note1},
already proved there and machine-checked in Lean. We therefore do not
\emph{avoid} Rado's direction; we \emph{re-prove} it from a lemma the paper
already owns, and \S\ref{sec:honest} says so plainly, replacing the incorrect
claim of independence in item 2 of the Gaps section of~\cite{note1}.
\end{abstract}

\section{What was missing, and what this note supplies}\label{sec:intro}

Theorem 7.1 of~\cite{note1} states, for $\lambda/\mu\in\cyl^{n,m}$, $\ell\ge 1$
and $\mathcal{T}_\ell(\lambda/\mu)\ne\emptyset$,
\begin{equation}\label{eq:thm71}
  \supp\bigl(s^{c}_{\lambda/\mu}(x_1,\dots,x_\ell)\bigr)\;=\;W_\ell
  \;=\;\{\alpha\in\N^{\ell}:\srt(\alpha)\trianglelefteq\lhat\}
  \;=\;\Perm_{\lhat}\cap\Z^{\ell},
\end{equation}
and concludes that $s^{c}_{\lambda/\mu}(x_1,\dots,x_\ell)$ is M-convex, has
saturated Newton polytope (SNP), and satisfies
$\Newton(s^{c}_{\lambda/\mu})=\Perm_{\lhat}$.

The second equality in~\eqref{eq:thm71} is proved there in full. The third is
not: the proof of Theorem 7.1 says only that it ``is the (elementary) statement
that the lattice points of the $\lhat$-permutahedron are exactly the $\alpha$
with $\srt(\alpha)\trianglelefteq\lhat$''. That statement is Rado's theorem
\cite{Rado}, and item~2 of the Gaps section of~\cite{note1} records it as
unproved. The last sentence of the proof then derives \emph{both} SNP and the
Newton identification from it, in the single clause ``SNP and the
identification of the Newton polytope follow since
$\Perm_{\lhat}=\conv(W_\ell)$''.

Two things are wrong with the way that was filed, and both matter more than the
missing lines themselves.

\begin{enumerate}
\item \textbf{The Gaps entry calls the missing statement ``cosmetic''.} Its exact
words are that it ``is used only for the cosmetic identification of the Newton
polytope''. But $\Newton(s^{c}_{\lambda/\mu})=\Perm_{\lhat}$ is one of the three
advertised conclusions of Theorem 7.1 and appears in the abstract of~\cite{note1}.
A conclusion a reader is invited to quote is not cosmetic. Grading it so is what
guaranteed that nobody --- including me --- ever asked whether the one-clause
derivation was valid.
\item \textbf{The same Gaps entry asserts that SNP does not depend on the missing
statement, and the text of~\cite{note1} does not realise that independence.}
The entry says ``M-convexity, SNP and the description of the support do not
depend on it''. The first and third are correct as written. The claim about SNP
is \emph{true} --- \S\ref{sec:snp} below gives a Rado-free proof --- but the
proof of Theorem 7.1 as printed derives SNP from
$\Perm_{\lhat}=\conv(W_\ell)$, i.e.\ exactly from the missing statement. So the
independence was asserted in the Gaps section and not enacted in the proof. This
is the precise content of the reviewer's remark that SNP holds ``only through a
line you do not write''.
\end{enumerate}

This note writes the lines. \S\ref{sec:JJ} isolates the sorted/subset identity.
\S\ref{sec:snp} proves SNP from it by convexity alone. \S\ref{sec:newton} proves
$\conv(W_\ell)=\Perm_{\lhat}$. \S\ref{sec:honest} states what has and has not
been achieved regarding independence from Rado. \S\ref{sec:verification} records
the computations. \S\ref{sec:gaps} records what remains.

\subsection*{Notation}
Throughout, $\lambda/\mu\in\cyl^{n,m}$, $\ell\ge 1$,
$\mathcal{T}_\ell(\lambda/\mu)\ne\emptyset$, and $d=|\lambda/\mu|$. For
$\alpha\in\Z^{\ell}$ and $S\subseteq[\ell]$ put $\alpha(S)=\sum_{i\in S}\alpha_i$,
and write $e_1,\dots,e_\ell$ for the standard basis. $\srt(\alpha)$ is the
partition obtained by sorting the positive entries of $\alpha$ decreasingly;
when convenient we regard a partition with at most $\ell$ parts as an element of
$\N^{\ell}$ by padding with zeros. Dominance $\sigma\trianglelefteq\nu$ is
defined for partitions \emph{of the same size} and means
$\sigma_1+\dots+\sigma_r\le\nu_1+\dots+\nu_r$ for all $r$. $\lhat$ is the sorted
greedy weight of Definition 5.1 of~\cite{note1},
$\Lambda_r=\lhat_1+\dots+\lhat_r$ for $0\le r\le\ell$, and
$\Perm_{\lhat}=\conv\{w\cdot\lhat : w\in S_\ell\}\subseteq\R^{\ell}$, where
$S_\ell$ acts by permuting coordinates. Finally
$W_\ell=W_\ell(\lambda/\mu)=\supp\bigl(s^{c}_{\lambda/\mu}(x_1,\dots,x_\ell)\bigr)$.

We use the following two results of~\cite{note1} and nothing else from it.

\begin{description}
\item[(A)] \textbf{Theorem 7.1, second equality:}
  $W_\ell=\{\alpha\in\N^{\ell}:\srt(\alpha)\trianglelefteq\lhat\}$. Proved
  in~\cite{note1} from Proposition 4.2 (the weight set is a dominance ideal),
  Proposition 5.5 (the greedy weight is the dominance maximum) and Corollary 3.4
  ($S_\ell$-stability). None of these consumes Rado's theorem, and the proof of
  the second equality does not use the third.
\item[(B)] \textbf{Lemma 4.1 (Robin Hood step):} if $\nu\rhd\sigma$ are
  partitions of $d$ then there are $a<b$ with $\nu_a\ge\nu_b+2$ such that
  $\tau=\nu-e_a+e_b$ is a partition and $\nu\rhd\tau\trianglerighteq\sigma$.
  Machine-checked: \texttt{RobinHood.robin\_hood\_step}, sorry-free, in
  \texttt{TworowD4Kernel/RobinHood.lean}.
\end{description}

Also from~\cite{note1}: Corollary 3.4, that $W_\ell$ is $S_\ell$-stable.

\section{The sorted and subset descriptions agree}\label{sec:JJ}

In the build of~\cite{note1} that this note repairs, the first sentence of the proof of the
polymatroid proposition (\texttt{prop:perm-mconvex}) asserted an
equality of two descriptions of the same set, justified by the clause ``because
$\max_{|S|=r}\alpha(S)$ is the sum of the $r$ largest entries of $\alpha$''. I
promote it here to a numbered lemma with a written proof, for two reasons. It is
the hypothesis of the SNP argument of \S\ref{sec:snp}; and, when this note was
written, it was the one step of that proposition \emph{not} formalised in Lean, the
formalisation working throughout in the subset form (\texttt{rem:mconvex-erratum}
of~\cite{note1}). \textbf{That second reason is no longer current and the correction
belongs here, not only in the Gaps section:} the lemma was transported the same evening as
\texttt{SortedBridge.insupp\_iff\_sorted} (commits \texttt{cc0573b}, \texttt{85054bb}),
sorry-free, and \S\ref{sec:honest} item 2 of the Gaps list records this. What is
\emph{not} formalised, and is unaffected by that closure, is Murota's \emph{symmetric}
exchange axiom --- not this lemma. The same bridge is
load-bearing in two places, and until now it was written out in neither.

\begin{lem}\label{lem:JJ}
Let $\lhat$ be a partition of $d$ with at most $\ell$ parts and
$\Lambda_r=\lhat_1+\dots+\lhat_r$. Then
\[
  \underbrace{\{\alpha\in\N^{\ell}:\srt(\alpha)\trianglelefteq\lhat\}}_{=:J}
  \;=\;
  \underbrace{\{\alpha\in\N^{\ell}:\alpha([\ell])=d,\
   \alpha(S)\le\Lambda_{|S|}\ \ \forall S\subseteq[\ell]\}}_{=:J'}.
\]
\end{lem}

\begin{proof}
First, for $\alpha\in\N^{\ell}$ and $0\le r\le\ell$,
\begin{equation}\label{eq:maxsum}
  \max_{S\subseteq[\ell],\,|S|=r}\alpha(S)\;=\;\srt(\alpha)_1+\dots+\srt(\alpha)_r,
\end{equation}
with the convention $\srt(\alpha)_k=0$ for $k$ beyond the number of positive
entries of $\alpha$. Indeed, choosing $S$ to consist of $r$ indices carrying $r$
largest entries realises the right-hand side; and for any $S$ with $|S|=r$ the
sum $\alpha(S)$ is a sum of $r$ entries of $\alpha$, hence at most the sum of the
$r$ largest. (The entries not among the $r$ largest may be $0$; this is why the
padding convention is the right one, and it is exactly where a partition with
fewer than $\ell$ parts is accommodated.)

$J\subseteq J'$. Let $\alpha\in J$, so $\srt(\alpha)\trianglelefteq\lhat$. Since
dominance is defined only between partitions of the same size, this \emph{includes}
the assertion $|\srt(\alpha)|=|\lhat|$, i.e.\ $\alpha([\ell])=d$. For
$S\subseteq[\ell]$ with $|S|=r$, \eqref{eq:maxsum} and dominance give
\[
  \alpha(S)\;\le\;\srt(\alpha)_1+\dots+\srt(\alpha)_r\;\le\;\lhat_1+\dots+\lhat_r
  \;=\;\Lambda_r\;=\;\Lambda_{|S|}.
\]

$J'\subseteq J$. Let $\alpha\in J'$. Then $|\srt(\alpha)|=\alpha([\ell])=d=|\lhat|$,
so the two are partitions of the same size and dominance between them is
meaningful. For $1\le r\le\ell$ choose, by \eqref{eq:maxsum}, a set $S$ with
$|S|=r$ and $\alpha(S)=\srt(\alpha)_1+\dots+\srt(\alpha)_r$; then
$\srt(\alpha)_1+\dots+\srt(\alpha)_r=\alpha(S)\le\Lambda_{|S|}=\Lambda_r$. For
$r>\ell$ both partial sums equal $d$, since $\srt(\alpha)$ and $\lhat$ each have
at most $\ell$ parts. Hence $\srt(\alpha)\trianglelefteq\lhat$.
\end{proof}

\begin{rem}\label{rem:equalsize}
The proof is short but it is not the sentence it replaces. What the one-sentence
justification omits is the role of the \emph{equal-size} convention: the
condition $\alpha([\ell])=d$ appears explicitly in $J'$ and only implicitly in
$J$, hidden inside the meaning of $\trianglelefteq$. Both inclusions above use
it, in opposite directions. Dropping it makes $J'\subseteq J$ false --- e.g.\
$\ell=2$, $\lhat=(2,1)$, $\alpha=(1,1)$ satisfies every inequality
$\alpha(S)\le\Lambda_{|S|}$ but $|\alpha|=2\ne 3$, and $(1,1)$ is not comparable
to $(2,1)$ in dominance at all.
\end{rem}

\section{Saturated Newton polytope}\label{sec:snp}

\begin{thm}[SNP]\label{thm:snp}
$\conv(W_\ell)\cap\Z^{\ell}=W_\ell$. Equivalently,
$s^{c}_{\lambda/\mu}(x_1,\dots,x_\ell)$ has saturated Newton polytope.
\end{thm}

\begin{proof}
Define
\[
  Q\;=\;\bigl\{x\in\R^{\ell}\;:\;x_i\ge 0\ (1\le i\le\ell),\quad x([\ell])=d,
  \quad x(S)\le\Lambda_{|S|}\ \ \forall S\subseteq[\ell]\bigr\}.
\]
$Q$ is an intersection of finitely many closed half-spaces with one hyperplane,
hence convex. By construction
\begin{equation}\label{eq:QZ}
  Q\cap\Z^{\ell}
  =\{\alpha\in\Z^{\ell}:\alpha_i\ge 0,\ \alpha([\ell])=d,\ \alpha(S)\le\Lambda_{|S|}\}
  =J',
\end{equation}
the middle set being exactly $J'$ because $\alpha\in\Z^{\ell}$ with all
$\alpha_i\ge0$ means $\alpha\in\N^{\ell}$.

By (A) and Lemma~\ref{lem:JJ}, $W_\ell=J=J'\subseteq Q$. Since $Q$ is convex and
$\conv(W_\ell)$ is the smallest convex set containing $W_\ell$, we get
$\conv(W_\ell)\subseteq Q$. Intersecting with $\Z^{\ell}$ and using
\eqref{eq:QZ},
\[
  \conv(W_\ell)\cap\Z^{\ell}\;\subseteq\;Q\cap\Z^{\ell}\;=\;J'\;=\;W_\ell .
\]
The reverse inclusion $W_\ell\subseteq\conv(W_\ell)\cap\Z^{\ell}$ is immediate
from $W_\ell\subseteq\N^{\ell}\subseteq\Z^{\ell}$. Hence equality.
\end{proof}

\begin{rem}\label{rem:snpfree}
Theorem~\ref{thm:snp} uses (A), Lemma~\ref{lem:JJ} and the convexity of a
half-space intersection. It does \emph{not} use $\conv(W_\ell)=\Perm_{\lhat}$,
does not use $\Perm_{\lhat}$ at all, and does not use Rado's theorem. This is the
sense in which the Gaps entry of~\cite{note1} was right to say SNP is
independent of item~2 --- but the route is the one above, not the one printed in
the proof of Theorem 7.1, which passes through $\conv(W_\ell)=\Perm_{\lhat}$.

Note also what is \emph{not} claimed: nothing above says $Q=\Perm_{\lhat}$. Only
$\conv(W_\ell)\subseteq Q$ is used, and that inclusion is all a lattice-point
argument needs. The equality $Q=\Perm_{\lhat}$ is true, but it is a statement
about real points and is strictly stronger than anything proved here; we neither
use nor prove it. (It follows from Theorem~\ref{thm:newton} together with
$\Perm_{\lhat}\subseteq Q$ and the fact that $Q$ is the convex hull of its
vertices, which are integral --- but that last fact is itself a theorem we have
not proved, so we stop.)
\end{rem}

\section{The Newton polytope}\label{sec:newton}

The content is the inclusion $W_\ell\subseteq\Perm_{\lhat}$. We prove it as an
instance of a statement about arbitrary convex $S_\ell$-stable sets, which is
where the argument is cleanest and which makes visible exactly which properties
of $\Perm_{\lhat}$ are used.

\begin{lem}[the segment step]\label{lem:segment}
Let $\nu\in\R^{\ell}$, let $a\ne b$ with $\nu_a>\nu_b$, and set
$g=\nu_a-\nu_b>0$ and $\tau=(a\,b)\cdot\nu$ (the transposition of coordinates
$a$ and $b$). Then
\[
  \nu-e_a+e_b\;=\;\Bigl(1-\tfrac{1}{g}\Bigr)\nu\;+\;\tfrac{1}{g}\,\tau ,
\]
and if $g\ge 1$ this is a convex combination of $\nu$ and $\tau$.
\end{lem}

\begin{proof}
$\tau$ agrees with $\nu$ outside $\{a,b\}$, and $\tau_a=\nu_b$, $\tau_b=\nu_a$,
so $\tau-\nu=(\nu_b-\nu_a)e_a+(\nu_a-\nu_b)e_b=g(e_b-e_a)$. Hence
\[
  \nu+\tfrac{1}{g}(\tau-\nu)=\nu+\tfrac{1}{g}\cdot g\,(e_b-e_a)=\nu-e_a+e_b,
\]
which is the displayed identity. If $g\ge 1$ then $\tfrac{1}{g}\in(0,1]$, so the
coefficients $1-\tfrac1g$ and $\tfrac1g$ are non-negative and sum to $1$.
\end{proof}

\begin{rem}\label{rem:gap1}
Lemma~\ref{lem:segment} needs only $g\ge 1$, whereas the Robin Hood step (B)
supplies $g=\nu_a-\nu_b\ge 2$. The stronger hypothesis is genuinely unused here:
asking which line of the proof of Lemma~\ref{lem:segment} breaks if $g=1$, the
answer is none --- one gets $\tfrac1g=1$ and the combination degenerates to
$\nu-e_a+e_b=\tau$, which is correct, since $\nu_a=\nu_b+1$ forces
$\nu-e_a+e_b=(a\,b)\cdot\nu$. The inequality $\nu_a\ge\nu_b+2$ is needed in the
proof of (B) itself, to make $\tau$ a partition \emph{strictly} below $\nu$; it
is not needed downstream. I record this because the reverse mistake --- carrying a
hypothesis into a lemma that does not consume it --- is how a hypothesis comes to
look like a constraint on the object.
\end{rem}

\begin{prop}\label{prop:minimal}
Let $K\subseteq\R^{\ell}$ be convex and stable under the coordinate action of
$S_\ell$, and let $\lhat\in K$ be a partition of $d$ with at most $\ell$ parts.
Then
\[
  \{\alpha\in\N^{\ell}:\srt(\alpha)\trianglelefteq\lhat\}\;\subseteq\;K .
\]
\end{prop}

\begin{proof}
\emph{Step 1: reduction to partitions.} Let $\alpha\in\N^{\ell}$ with
$\srt(\alpha)\trianglelefteq\lhat$ and put $\sigma=\srt(\alpha)\in\N^{\ell}$
(padded). Then $\alpha=w\cdot\sigma$ for some $w\in S_\ell$, so by
$S_\ell$-stability of $K$ it suffices to prove $\sigma\in K$. So let $\sigma$ be
a partition of $d$ with at most $\ell$ parts and $\sigma\trianglelefteq\lhat$.

\emph{Step 2: the chain.} Define $\nu^{(0)}=\lhat$ and, recursively, as long as
$\nu^{(t)}\ne\sigma$: we will maintain the invariant that $\nu^{(t)}$ is a
partition of $d$ with at most $\ell$ parts and
$\nu^{(t)}\trianglerighteq\sigma$. Given the invariant and $\nu^{(t)}\ne\sigma$
we have $\nu^{(t)}\rhd\sigma$, so (B) applies and supplies $a_t<b_t$ with
$\nu^{(t)}_{a_t}\ge\nu^{(t)}_{b_t}+2$ such that
\[
  \nu^{(t+1)}\;:=\;\nu^{(t)}-e_{a_t}+e_{b_t}
\]
is a partition with $\nu^{(t)}\rhd\nu^{(t+1)}\trianglerighteq\sigma$. The
invariant is preserved: $\nu^{(t+1)}$ is a partition of $d$ dominating $\sigma$,
and it has at most $\ell$ parts because $b_t\le\ell$. (For the latter: in the
proof of (B) the index $b$ satisfies $b\le j$, where $j$ is chosen with
$\nu^{(t)}_j<\sigma_j$; since $\sigma$ has at most $\ell$ parts, $\sigma_j>0$
forces $j\le\ell$.) Also $\nu^{(t+1)}\in\N^{\ell}$, since
$\nu^{(t+1)}_{a_t}=\nu^{(t)}_{a_t}-1\ge\nu^{(t)}_{b_t}+1\ge 1$.

\emph{Step 3: termination.} Write $N^{(t)}_r=\nu^{(t)}_1+\dots+\nu^{(t)}_r$ and
$S_r=\sigma_1+\dots+\sigma_r$, and set
\[
  \Phi(t)\;=\;\sum_{r=1}^{\ell}\bigl(N^{(t)}_r-S_r\bigr).
\]
Each summand is $\ge 0$ because $\nu^{(t)}\trianglerighteq\sigma$, so
$\Phi(t)\in\N$. The partial sums of $\nu^{(t+1)}=\nu^{(t)}-e_{a_t}+e_{b_t}$ are
$N^{(t)}_r-[a_t\le r<b_t]$, so
\[
  \Phi(t+1)\;=\;\Phi(t)-\#\{r: a_t\le r<b_t\}\;=\;\Phi(t)-(b_t-a_t)\;\le\;\Phi(t)-1,
\]
using $a_t<b_t\le\ell$, so that all of $a_t,\dots,b_t-1$ lie in $[1,\ell]$.
Thus $\Phi$ is a strictly decreasing sequence of non-negative integers, and the
recursion halts after at most $\Phi(0)$ steps. It halts only when
$\nu^{(k)}=\sigma$ for some $k\le\Phi(0)$.

\emph{Step 4: the chain stays in $K$.} We show $\nu^{(t)}\in K$ by induction on
$t$. For $t=0$, $\nu^{(0)}=\lhat\in K$ by hypothesis. Suppose $\nu:=\nu^{(t)}\in
K$ and write $a=a_t$, $b=b_t$, $g=\nu_a-\nu_b\ge 2$. Since $K$ is
$S_\ell$-stable, $(a\,b)\cdot\nu\in K$. By Lemma~\ref{lem:segment},
$\nu^{(t+1)}=\nu-e_a+e_b$ is a convex combination of $\nu$ and
$(a\,b)\cdot\nu$, both in $K$; as $K$ is convex, $\nu^{(t+1)}\in K$.

By Step 3 the chain reaches $\sigma$, so $\sigma=\nu^{(k)}\in K$, and by Step 1
$\alpha\in K$.
\end{proof}

\begin{cor}\label{cor:minimalpoly}
$\Perm_{\lhat}$ is the smallest convex $S_\ell$-stable subset of $\R^{\ell}$
containing $\lhat$, and it contains $\{\alpha\in\N^{\ell}:\srt(\alpha)\trianglelefteq\lhat\}$.
\end{cor}

\begin{proof}
$\Perm_{\lhat}=\conv(S_\ell\cdot\lhat)$ is convex, and it is $S_\ell$-stable
because $w\cdot\conv(S_\ell\cdot\lhat)=\conv(w S_\ell\cdot\lhat)=\conv(S_\ell\cdot\lhat)$
for $w\in S_\ell$ (the action is linear, so it commutes with taking convex
hulls). It contains $\lhat$. Conversely any convex $S_\ell$-stable $K\ni\lhat$
contains $S_\ell\cdot\lhat$ and hence its convex hull $\Perm_{\lhat}$. The second
assertion is Proposition~\ref{prop:minimal} with $K=\Perm_{\lhat}$.
\end{proof}

\begin{thm}[Newton polytope]\label{thm:newton}
$\conv(W_\ell)=\Perm_{\lhat}$; that is,
$\Newton\bigl(s^{c}_{\lambda/\mu}(x_1,\dots,x_\ell)\bigr)=\Perm_{\lhat}$.
Moreover $W_\ell=\Perm_{\lhat}\cap\Z^{\ell}$, which is the third equality
of~\eqref{eq:thm71}.
\end{thm}

\begin{proof}
$\Perm_{\lhat}\subseteq\conv(W_\ell)$: by (A), $\lhat\in W_\ell$ (indeed
$\srt(\lhat)=\lhat\trianglelefteq\lhat$), and $W_\ell$ is $S_\ell$-stable by
Corollary 3.4 of~\cite{note1}, so $S_\ell\cdot\lhat\subseteq W_\ell$ and
therefore $\Perm_{\lhat}=\conv(S_\ell\cdot\lhat)\subseteq\conv(W_\ell)$.

$\conv(W_\ell)\subseteq\Perm_{\lhat}$: by (A),
$W_\ell=\{\alpha\in\N^{\ell}:\srt(\alpha)\trianglelefteq\lhat\}$, which is
contained in $\Perm_{\lhat}$ by Corollary~\ref{cor:minimalpoly}. Taking convex
hulls of $W_\ell\subseteq\Perm_{\lhat}$ and using convexity of $\Perm_{\lhat}$
gives $\conv(W_\ell)\subseteq\Perm_{\lhat}$.

For the last assertion, $W_\ell\subseteq\Perm_{\lhat}\cap\Z^{\ell}$ is now
immediate, and conversely
$\Perm_{\lhat}\cap\Z^{\ell}=\conv(W_\ell)\cap\Z^{\ell}=W_\ell$ by
Theorem~\ref{thm:snp}.
\end{proof}

\begin{rem}
The two theorems are logically independent in the following sense.
Theorem~\ref{thm:snp} does not use Theorem~\ref{thm:newton}: it never mentions
$\Perm_{\lhat}$. Theorem~\ref{thm:newton} does use Theorem~\ref{thm:snp}, but
only for its final sentence, the lattice-point identity; the polytope identity
$\conv(W_\ell)=\Perm_{\lhat}$ is proved without it. So the printed one-clause
derivation in~\cite{note1}, which obtains SNP \emph{from}
$\conv(W_\ell)=\Perm_{\lhat}$, had the dependency the wrong way round: SNP is the
cheaper statement and should come first.
\end{rem}

\section{What is and is not independent of Rado}\label{sec:honest}

Item~2 of the Gaps section of~\cite{note1} reads, in part, that the missing
statement ``is Rado's theorem, and it is used only for the cosmetic
identification of the Newton polytope''. I replace it with the following, which
is what the present note actually establishes.

\begin{enumerate}
\item \textbf{SNP is genuinely independent of Rado.} Theorem~\ref{thm:snp}
consumes only (A), Lemma~\ref{lem:JJ} and the convexity of an intersection of
half-spaces. No permutahedron, no majorisation theorem.

\item \textbf{The Newton identification is not independent of Rado's content ---
it is a re-proof of it.} Proposition~\ref{prop:minimal} with
$K=\Perm_{\lhat}$ \emph{is} the non-trivial direction of Rado's theorem
\cite{Rado}, equivalently of the Hardy--Littlewood--P\'olya majorisation theorem
\cite[Thm.~46]{HLP} and of Birkhoff's theorem in the form usually quoted
alongside them: $\srt(\alpha)\trianglelefteq\lhat$ implies
$\alpha\in\conv(S_\ell\cdot\lhat)$. Moreover the proof given above is the
classical one: the Robin Hood step (B) is the classical transfer (or Muirhead,
or Dalton) step, and Step 4 is the classical observation that a transfer moves
along a segment between a point and its transposition.

So the honest sentence is: \emph{we re-prove Rado's direction from Lemma 4.1}.
It is not \emph{we avoid Rado}.

\item \textbf{Nevertheless the self-containedness claim of~\cite{note1}
survives, and is what was actually wanted.} What~\cite{note1} claims, and what
WZZ Problem 5.2 \cite{WZZ} asks for, is a proof that \emph{cites} nothing beyond
the tableau combinatorics --- in particular not Lam's Corollary 8.5 and not Lee's
Corollary 5. That is unaffected: Lemma 4.1 is proved in~\cite{note1} from
scratch for an entirely different purpose (Proposition 4.2, that the weight set
is a dominance ideal), and it turns out to drive the majorisation argument too.
The paper cites nothing; it simply happens to re-prove a classical theorem along
the way, and should say so rather than claim to have sidestepped it.

\item \textbf{The word ``cosmetic'' must go.}
$\Newton(s^{c}_{\lambda/\mu})=\Perm_{\lhat}$ is one of the three advertised
conclusions of Theorem 7.1, and it appears in the abstract. Filing a headline
conclusion as cosmetic is what kept the one-clause derivation from ever being
read.
\end{enumerate}

\section{Computational verification}\label{sec:verification}

All arithmetic below is exact (Python integers and \texttt{fractions.Fraction});
no floating point and no LP solver is used anywhere. Code:
\texttt{projects/proofs/code-q261/rado\_from\_robin\_hood.py} and
\texttt{projects/proofs/code-q261/real\_supports\_snp\_newton.py}.

\paragraph{V1 --- Lemma~\ref{lem:JJ}.}
$J$ and $J'$ computed independently, $J'$ by testing all $2^{\ell}$ subsets
rather than via \eqref{eq:maxsum}, so the check does not presuppose the lemma's
own justification. All $\lhat$ with $|\lhat|\le 9$ and all $\ell\le 5$:
$247$ pairs $(\lhat,\ell)$, $44245$ points $\alpha$, \textbf{0 disagreements}.

\paragraph{V2 --- the chain of Proposition~\ref{prop:minimal}.}
For every pair of partitions $\sigma\trianglelefteq\lhat$ of $d\le 9$ with at
most $\ell\le 4$ parts ($725$ pairs, $1174$ Robin Hood steps), the construction
of Step 2 was run with the three proof obligations asserted at every step:
$a_t<b_t$; the gap $\nu^{(t)}_{a_t}-\nu^{(t)}_{b_t}\ge 2$ (so
$1/g\in(0,\tfrac12]$); $\nu^{(t+1)}$ a partition in $\N^{\ell}$; and
$\Phi(t+1)=\Phi(t)-(b_t-a_t)<\Phi(t)$. The induction of Step 4 was then
\emph{unrolled} into an explicit convex combination: each $\nu^{(t)}$ is carried
as a dictionary from points of $S_\ell\cdot\lhat$ to rational coefficients, and
each step applies
$\nu^{(t+1)}=(1-\tfrac1g)\nu^{(t)}+\tfrac1g\,(a_t\,b_t)\nu^{(t)}$ termwise.
Checked for each pair: coefficients non-negative and summing to exactly $1$;
every point of the support an honest permutation of $\lhat$; and the resulting
combination equal to $\sigma$ exactly. \textbf{0 failures}; maximum support size
$9$.

This is the check that matters, and it is deliberately not a feasibility test. A
solver confirming $\sigma\in\Perm_{\lhat}$ would verify the \emph{conclusion} of
Proposition~\ref{prop:minimal} while saying nothing about whether \emph{this}
argument produces it. In particular it would not detect the third of the failure
modes an argument of this shape is prone to: terminating at dominance-maximal
elements rather than at permutations of $\lhat$, so that the combination
produced is of points of $J$ and not of vertices. V2 tests exactly that, by
requiring every support point to be a permutation of $\lhat$.

\paragraph{V3 --- $Q\cap\Z^{\ell}=J'$.}
Every lattice point of $Q$ enumerated for $|\lhat|\le 8$, $\ell\le 4$:
$128$ pairs, $3303$ lattice points, \textbf{0 outside $J$}.

\paragraph{V4 --- both theorems on the real supports.}
The weight sets $W_\ell$ were produced by an independent enumerator
(\texttt{peers/rick/proofs/2026-09-25-clio288-greedy-check.py}, by Rick), which
builds chains from the raw horizontal-strip condition and never uses Lemma 3.1
of~\cite{note1} or the greedy chain. Deliberately, the bounding inequalities
were read off $W_\ell$ itself,
\[
  M_r\;:=\;\max_{\alpha\in W_\ell}\ (\text{sum of the $r$ largest entries of }\alpha),
\]
with no reference to $\lhat$, so that the SNP test does not route through (A).
Checked for each instance: (i) $\{\alpha\in\N^{\ell}:|\alpha|=d,\ \alpha(S)\le
M_{|S|}\ \forall S\}=W_\ell$ (Theorem~\ref{thm:snp} on the real object); (ii)
$M_r=\Lambda_r$ for all $r$ (a consequence, checked rather than assumed); (iii)
for every $\alpha\in W_\ell$, the V2 certificate --- an exact rational convex
combination of permutations of $\lhat$ equal to $\alpha$; (iv) every permutation
of $\lhat$ lies in $W_\ell$, which is the easy inclusion of
Theorem~\ref{thm:newton}.

Over $n\le 7$, $d\le 8$, $\ell\le 5$: \textbf{9228 instances
$(n,m,\mu,\lambda,\ell)$ with $W_\ell\ne\emptyset$, 268319 lattice points,
268319 exact certificates, 0 failures} in all four checks.

\paragraph{What these checks cannot see.} They are all instances of finite
$(n,d,\ell)$. Lemma~\ref{lem:JJ}, Lemma~\ref{lem:segment} and
Proposition~\ref{prop:minimal} are proved above for all parameters, and the
computations are calibration, not evidence. Nothing here tests (A) or (B) ---
(A) is proved in~\cite{note1} and (B) is additionally machine-checked.

\section{Gaps}\label{sec:gaps}

\begin{enumerate}
\item \textbf{$Q=\Perm_{\lhat}$ is not proved.} Remark~\ref{rem:snpfree}. We
prove $\conv(W_\ell)\subseteq Q$ and $\conv(W_\ell)=\Perm_{\lhat}$, hence
$\Perm_{\lhat}\subseteq Q$; the reverse inclusion --- every \emph{real} point of
$Q$ lies in $\Perm_{\lhat}$ --- is the real-coefficient form of Rado's theorem
and is not needed for anything in~\cite{note1}. It is not claimed anywhere and
should not be quoted from this note.
\item \textbf{Lemma~\ref{lem:JJ} is formalised --- this gap closed on 25 September
2026, cycle 2, a few hours after this note was written.} As printed here the item
read ``still not formalised'', and said that \texttt{polymatroid-exchange} would
stay below \texttt{lean-verified} until Lemma~\ref{lem:JJ} was transported. It was
transported that evening: \texttt{SortedBridge.insupp\_iff\_sorted} and
\texttt{SortedBridge.sorted\_exchange} in
\texttt{TworowD4Kernel/SortedSubsetBridge.lean} (commits \texttt{cc0573b},
\texttt{85054bb}), sorry-free. The prediction was correct and the sentence went
stale by being acted on; it is recorded here rather than deleted, because a caveat
is a claim about the state of another file and its own success is what refutes it.
What remains outside the Lean record is Murota's \emph{symmetric} exchange axiom,
not this lemma.
\item \textbf{Proposition~\ref{prop:minimal} is classical and is not claimed as
new.} \S\ref{sec:honest}, item 2. Its only novelty here is that its engine was
already available in~\cite{note1}, so that the paper needs no citation.
\end{enumerate}

\begin{thebibliography}{9}
\bibitem{note1} Clio, \emph{A combinatorial proof of the M-convexity of
cylindric skew Schur polynomials}, \texttt{projects/proofs/2026-09-20-c1-cylindric-M-convexity.tex}.
Cited by internal label, not by number, because folding the present results in
as \S8 of that note renumbers part of \S6. For the build this note is committed alongside
(\texttt{clio-vega/proofs@a34a66d}) the correspondence is: \texttt{lem:hlp} = Lemma 4.1,
\texttt{prop:ideal} = Proposition 4.2, \texttt{cor:bk} = Corollary 3.4, \texttt{lem:greedy} =
Lemma 5.2, \texttt{prop:max} = Proposition 5.5, \texttt{lem:tight} = Lemma 6.1, \texttt{lem:JJ} =
Lemma 6.2 (new), \texttt{prop:perm-mconvex} = Proposition 6.3 (\textbf{was 6.2 before the
insertion of \texttt{lem:JJ}}), \texttt{rem:mconvex-erratum} = Remark 6.4 (\textbf{was 6.3}),
\texttt{thm:main} = Theorem 7.1 (unchanged). All numbers read from that build's \texttt{.aux}
file, not from recall. Note that \S8 of~\cite{note1} contains the same mathematics as
\S\S\ref{sec:JJ}--\ref{sec:newton} here, so its Lemma 8.2, Proposition 8.3 and Theorems 8.1, 8.4
are this note's Lemma~\ref{lem:segment}, Proposition~\ref{prop:minimal},
Theorem~\ref{thm:snp} and Theorem~\ref{thm:newton}; and its Lemma 6.2 is this note's
Lemma~\ref{lem:JJ}.
\bibitem{Rado} R.~Rado, \emph{An inequality}, J.\ London Math.\ Soc.\ \textbf{27}
(1952), 1--6. (Cited for attribution of the majorisation statement re-proved in
\S\ref{sec:newton}. Not read at source; the attribution is the standard one and
is used here only to name a classical theorem, not to import a step. Extraction
level: reference-only.)
\bibitem{HLP} G.~H.~Hardy, J.~E.~Littlewood, G.~P\'olya, \emph{Inequalities},
2nd ed., Cambridge, 1952, Theorem 46. (Same status as \cite{Rado}:
attribution only.)
\bibitem{WZZ} Wang, Zhang, Zhang, \emph{Newton polytopes and M-convexity of dual
$k$-Schur polynomials}, arXiv:2401.14632. (Read at LaTeX source; Problems
5.1--5.2 at lines 938 and 945.)
\end{thebibliography}

\end{document}
